\begin{abstract}
Household food waste motivates the development of tools that connect available ingredients with practical meal planning. However, generated recipes require verification against ingredient quantities, storage information, and user constraints before they can support actionable decisions. This study proposes Food Guardian, an inventory-aware artificial intelligence agent that combines natural-language interaction with database retrieval, food-storage rules, and usage analysis. A large language model coordinates tool calls and generates candidate meals, while deterministic modules check ingredient availability, enforce defined constraints, and prioritize eligible ingredients by relevant dates. Missing storage information triggers clarification, and inventory changes require user confirmation. The prototype architecture is designed using Python, SQLite, Pandas, and a source-documented rule base. Evaluation employs 20--30 ingredient types, 20--30 recipes, and 30 benchmark scenarios covering sufficient inventory, shortages, date-related conditions, incomplete information, and preference conflicts. Experimental comparisons contrast a generic language model ($B_0$) and an inventory-conditioned model ($B_1$) against the full Food Guardian pipeline ($P$). Assessment measures recipe feasibility, inventory matching, date-reminder accuracy, rule violations, clarification behavior, and feasible-solution coverage. The expected contributions are an explainable agent prototype, traceable decision records, and a reproducible evaluation procedure for studying the integration of language-model planning with explicit inventory constraints.
\end{abstract}

\begin{IEEEkeywords}
Artificial intelligence agents, household food waste, inventory management, recipe recommendation, food-storage rules, tool calling.
\end{IEEEkeywords}
